% =============================================================================
% si.tex — Supplementary Information (draft v0.3, 2026-09-03)
% Data description, parameter tables, solver settings, demand data pack
% (v0.3: SI tables renumbered S1..S4; Table S1/S2/S3 bodies added 2026-09-03)
% (v0.4 2026-09-04: SI figures renumbered S-Figure; tornado moved here per Ma-style)
% =============================================================================
\renewcommand{\thetable}{S\arabic{table}}
\setcounter{table}{0}
\renewcommand{\thefigure}{S\arabic{figure}}
\setcounter{figure}{0}

\section{Data sources and processing}

\subsection{Biomass availability}

County-level feedstock availability (dry tonnes/yr) comes from the DOE 2023
Billion-Ton Report (BT23) county-level downloads
(\texttt{bioenergykdf.ornl.gov}), agricultural and forestry modules; the BT23
forestry module draws on USFS Forest Inventory and Analysis (FIA) data, which
is therefore the underlying forestry source cited in the main text. The
baseline uses the \emph{near-term} scenario; the \emph{mature-market medium}
scenario serves as sensitivity. The BT23 corn stover series is constructed
from grain-corn production via a 1:1 stover-to-grain ratio and is therefore
\emph{exclusive of silage corn acreage by construction} (silage is absent from
the POLYSYS crop list; the 2022 Census of Agriculture reports 0.81~M acres of
corn silage versus 3.23~M acres of grain corn harvested (20.1\% of corn
acreage), so erroneously including silage would inflate stover by roughly
25\%; the county-level grain/silage split is archived for a silage-inclusive
sensitivity). The download field \texttt{price\_offered} is the
farmgate bid threshold (2022 USD per dry short ton, ash-inclusive, excluding
transport); the near-term scenario represents approximately 2030 conditions
with residues only and a \$70/dry-ton price ceiling. The tonnages are
theoretical availability: field-level collection losses, residue-retention
requirements, and landowner participation are not modeled (a $-40\%$
availability sensitivity is reported in the main text). Feedstock mapping: corn
stover, miscanthus, switchgrass, willow, and poplar map directly from the
agricultural download; logging residues aggregate the five BT23 logging-residue
classes (hardwood lowland/upland, mixedwood, softwood natural/planted);
small-diameter trees aggregate the four SDT classes; mill residues aggregate
hardwood and softwood processing residues; other forest waste corresponds to
``forest waste human generated.'' Wheat, oats, barley, and soybean straw are
taken from USDA NASS 2022 Census of Agriculture county production data
converted to theoretical residue tonnage (NASS Census residue estimates,
archived in the project workbook sheet \texttt{02\_NASS\_Validation};
Table~S1 gives the county corn grain/silage acreage split),
because BT23 reports near-zero availability for these resources at \$70/t.
State totals: near-term 6.97~Mt dry/yr; mature-market medium 13.53~Mt dry/yr
(of which the nine modeled BT23 categories account for 12.54~Mt). BT23
\emph{waste} resources are characterized but not yet modeled: the full
Wisconsin wastes download (16 categories, 72 counties, five scenarios;
archived as \texttt{bt23\_wastes\_WI\_summary.csv}) gives 2.27~Mt dry/yr at
$\le$\$70/t in the near-term scenario, dominated by dairy manure
(1.13~Mt), clean urban wood (0.34~Mt), textiles (0.27~Mt), and yard waste
(0.25~Mt), with sludge, paper, plastics, and FOG at negligible $\le$\$70/t
availability; these form a co-location extension for Wisconsin's dairy and
paper economy (Section Limitations) and are an explicit model-boundary
statement, not an oversight of data.

\begin{table}[htbp]
\centering
\caption{\small USDA NASS 2022 Census of Agriculture, Wisconsin corn acreage split
(archived \texttt{corn\_grain\_silage\_county\_2022.csv}). The BT23 grain-corn
stover series excludes silage by construction; a silage-inclusive stover
estimate would add roughly 25\%.}
\label{tab:s1}
\begin{tabular}{lcc}
\toprule
 & acres (M) & share of corn acreage\\
\midrule
Grain corn harvested & 3.23 & 79.9\%\\
Corn silage & 0.81 & 20.1\%\\
\bottomrule
\end{tabular}
\end{table}

\subsection{Feedstock characterization}

Table~S2 gives, per feedstock: initial moisture content (12--45\%), drying
yield to 10\% moisture $\eta=(100-\mathrm{MC})/90$, farmgate prices (\$40--60/dry
t, wet-adjusted), 300$^\circ$C biochar yield (0.33--0.38 t/t), carbon content
(0.41--0.50), and molar H/C (1.00--1.40 at 300$^\circ$C, 0.40--0.52 at
500$^\circ$C). Moisture contents follow the
feedstock-collection literature compiled in Roberts et al.\ \cite{roberts2010}
(corn stover 15--30\%, woody biomass 45\%) and the forestry-drying review of
Fagernas et al.\ \cite{fagernas2010}; char carbon contents are cross-checked
against the feedstock-class means of Woolf et al.\ \cite{woolf2010} (SI Table
S4: cereals 59.9\%, herbaceous crops 56.8\%, wood 75.0\%, green waste 67.5\%).
Yields are expressed per tonne of 10\%-moisture dried
feedstock (the model's dry product); if oven-dry-basis literature yields are
preferred, all yield and per-dry GHG factors scale by 0.9.

\begin{table}[htbp]
\centering
\small
\caption{\small Feedstock characterization: moisture (MC), dry carbon content
$C_{\mathrm{dry}}$, biochar mass yields $y_{300}/y_{500}$, molar H/C at
300/500$^\circ$C, and farmgate price. H/C values follow the literature audit of
Section~A.3 (300$^\circ$C: Rafiq et al.\ 2016 and the Ippolito et al.\ 2020
meta-analysis band; 500$^\circ$C: Rafiq 2016, Enders 2012, Ippolito 2020): all
300$^\circ$C chars exceed the VM0044 H/C$\le$0.7 gate (torrefaction-like, never
credit-eligible), and all 500$^\circ$C chars satisfy it.}
\label{tab:s2}
\begin{tabular}{llcccccc cc}
\toprule
Feedstock & class & MC (\%) & $C_{\mathrm{dry}}$ & $y_{300}$ & $y_{500}$ & H/C$_{300}$ & H/C$_{500}$ & farmgate & Reference\\
 & & & & & & & & (\$/dry t) & \\
\midrule
Corn stover & ag & 17.5 & 0.45 & 0.38 & 0.300 & 1.40 & 0.49 & 60 & \cite{rafiq2016,ippolito2020}\\
Soybean straw & ag & 14.0 & 0.43 & 0.36 & 0.284 & 1.30 & 0.50 & 60 & \cite{ippolito2020}\\
Wheat straw & ag & 12.5 & 0.42 & 0.35 & 0.276 & 1.25 & 0.51 & 60 & \cite{ippolito2020}\\
Oats straw & ag & 13.0 & 0.41 & 0.34 & 0.268 & 1.25 & 0.52 & 60 & \cite{ippolito2020}\\
Barley straw & ag & 12.0 & 0.41 & 0.34 & 0.268 & 1.25 & 0.51 & 60 & \cite{ippolito2020}\\
Switchgrass & ag & 12.5 & 0.47 & 0.38 & 0.300 & 1.30 & 0.47 & 60 & \cite{ippolito2020}\\
Miscanthus & ag & 15.0 & 0.48 & 0.37 & 0.292 & 1.30 & 0.45 & 60 & \cite{ippolito2020}\\
Poplar & for & 45.0 & 0.50 & 0.35 & 0.276 & 1.10 & 0.40 & 60 & \cite{woolf2010,enders2012}\\
Willow & for & 45.0 & 0.50 & 0.35 & 0.276 & 1.10 & 0.40 & 60 & \cite{woolf2010,enders2012}\\
Logging residues & for & 42.5 & 0.49 & 0.33 & 0.261 & 1.05 & 0.42 & 40 & \cite{fagernas2010,enders2012}\\
Small-diameter trees & for & 45.0 & 0.49 & 0.33 & 0.261 & 1.05 & 0.41 & 60 & \cite{roberts2010,enders2012}\\
Mill residues & for & 35.0 & 0.50 & 0.34 & 0.268 & 1.00 & 0.40 & 50 & \cite{enders2012}\\
Other forest waste & for & 20.0 & 0.47 & 0.35 & 0.276 & 1.15 & 0.45 & 50 & \cite{ippolito2020}\\
\bottomrule
\end{tabular}
\end{table}

\subsection{Technology economics}

T1 dehydration: OPEX \$5.2--10.8/wet t (moisture-dependent), CAPEX
\$0.8/2.5/7.0\,M for 33/165/660\,kt/yr plants \cite{yue2014,gebreegziabher2014}.
T2 slow
pyrolysis: OPEX \$30/dry t, CAPEX \$9/21.5/51\,M for 60/198/660\,kt/yr
\cite{shabangu2014,roberts2010}. The 300$^\circ$C/500$^\circ$C yield and
carbon-content pairs
for corn stover (0.38/0.30~t/t; 0.45/0.55) are the advisor-confirmed
technology-table values (V9 workbook, sheet 02); per-feedstock 500$^\circ$C
values use the documented scaling
$y_{500}=0.789\,y_{300}$, $C_{500}=1.222\,C_{300}$. Source audit: Roberts et
al.\ \cite{roberts2010} is a 450$^\circ$C slow-pyrolysis LCA and contains no
300$^\circ$C or 500$^\circ$C case; it supports $y_{500}=0.30$ indirectly, via
its char-yield range of 28.8--33.0~wt\% (citing Antal and Gr{\o}nli 2003,
\emph{Ind.\ Eng.\ Chem.\ Res.}\ 42:1619), it matches the Rafiq et al.\ corn-stover
500$^\circ$C yield of 29.2\% \cite{rafiq2016}, and it supports
$\rho_{500}=0.80$ through its 80\% stable-carbon assumption. Woolf et
al.\ \cite{woolf2010} do not tabulate temperature-series yields or H/C
ratios. The H/C values are taken from the temperature-series literature
(Table~S2): 300$^\circ$C chars have H/C $=1.00$--$1.40$ (torrefaction-like;
Rafiq et al.\ \cite{rafiq2016} measure 1.41 for corn stover at 300$^\circ$C,
and the Ippolito et al.\ meta-analysis \cite{ippolito2020} places sub-400$^\circ$C
chars in the ``partially thermochemically altered'' range), so
\emph{no} 300$^\circ$C biochar meets the VM0044 H/C$\le$0.7 gate, and its
credit rate is zero throughout. The 500$^\circ$C chars have H/C
$=0.40$--$0.52$ (Rafiq et al.\ \cite{rafiq2016}: corn stover 0.49, benchmarked
against Enders et al.\ \cite{enders2012}; the Ippolito et al.\ meta-analysis
500--599$^\circ$C band mean is $\approx$0.54), comfortably below the gate, so
all 500$^\circ$C biochar is credit-eligible. A $\pm$0.1 H/C perturbation
changes no eligibility outcome (Section~A.7); the first flip of a 500$^\circ$C
char requires $+0.25$, which is the stress test reported there. The
$C_{500}=0.55$ value is conservative relative to the published 500$^\circ$C
carbon contents of 62.5\% (Ippolito et al.\ \cite{ippolito2020} band mean)
and 64.5\% (Rafiq et al.\ \cite{rafiq2016}), so 500$^\circ$C credit rates are
conservative lower bounds.

\subsection{Greenhouse gas factors}

Transport 0.0902~kg CO$_2$e/t-km (32-t truck) \cite{wernet2016}; collection and
baling
35~kg/t dry \cite{yue2014}; drying 25~kg/t H$_2$O
\cite{gebreegziabher2014,fagernas2010}; slow pyrolysis 150~kg/t dry net of a
syngas electricity credit at 50--65~kWh/t and 0.57~kg/kWh Wisconsin grid
intensity (Roberts et al.~\cite{roberts2010}); (the pyrolysis term is the dominant process
emission and carries the largest parameter uncertainty). Carbon removal is
defined as a positive quantity,
$r^{\mathrm{seq}}_t=\rho_t\,C_t\,(44/12)\times 10^{3}$ kg CO$_2$e/t biochar,
with permanence
$\rho_{300}=0.65$, $\rho_{500}=0.80$ per VM0044 v1.2 Table 3 (IPCC 2019
Table 4AP.2: 0.65 for 350--450$^\circ$C, 0.80 for 450--600$^\circ$C, 0.89 for
$>$600$^\circ$C; the 300$^\circ$C value uses the lowest tabulated band,
350--450$^\circ$C, and the permanence of sub-350$^\circ$C chars may be lower,
so $\rho_{300}$ should be read as an upper-bound parameter); N$_2$O suppression
$52$~kg/t \cite{cayuela2014};
decomposition baseline $b_f=0.9\,C_f^{\mathrm{dry}}(44/12)\times 10^{3}$
kg/t dry. Carbon credits are intangible certificates: they carry zero transport
cost and are \emph{not} assigned transport emissions (the certificates are
settled at the producing county).

\subsection{Demand data pack}

The three-segment biochar demand curve is calibrated to: (i) production-side
costs of \$150--400/t \cite{roberts2010,shabangu2014} plus the removal-credit
value of \$100--250/tCO$_2$e \cite{cdrfyi2024,verraVM0044}, which places the
durable-char tiers above the bulk soil-amendment tier; (ii) the
soil-amendment, industrial-media, and carbon-removal end uses that structure
the biochar market, with wholesale prices in
recent market surveys spanning more than an order of magnitude
\cite{ibi2025}; (iii)
Wisconsin farmland of 14.2~M acres
with application rates of 5--20~t/ha, within the range of field studies
reviewed by Jeffery et al.~\cite{jeffery2011}, and adoption below 100\%, so
that the
4.16~Mt aggregate of the previous model version is treated as an optimistic
peak; and (iv) a non-soil sink valued at \$30/t (construction/filtration media)
with a landfill disposal sensitivity of $-\$45$ to $-\$84$/t. Segment
capacities are allocated across counties in proportion to farmland area; the
CDR-premium segment is interpreted as access to national corporate buyers
through credit markets rather than local Wisconsin demand, and the farmland
proportionality is a spatial proxy (county-specific adoption data are not
available). Because every county carries demand at the same bid, the binding
demand constraint is aggregate and the equilibrium spatial pattern is
feedstock-driven; the concentration sensitivity below tests this directly.
The 0.8~Mt premium capacity is a national-market allocation:
current U.S. biochar purchases reported by CDR.fyi are in the hundreds of
thousands of tonnes per year \cite{cdrfyi2024}, so the segment assumes the
several-fold market growth projected for biochar credits rather than today's
volumes. The decomposition fraction $0.9$ is an upper bound (immediate,
complete decay); a sensitivity to $0.7$/$0.5$ scales the Paradigm~A credit
rates proportionally and is flagged in the Limitations section.

\emph{Demand-geography sensitivity.} The baseline allocates each demand
segment across counties in proportion to farmland area ($w_n \propto
\mathrm{area}_n$). To test whether the equilibrium depends on this spatial
proxy, we re-solve the Step-1 MIP and Step-2 LP with demand concentrated as
$w_n \propto \mathrm{area}_n^{\alpha}$ for $\alpha \in \{3, 10\}$; at
$\alpha=10$ the three largest counties carry 71\% of the H-segment capacity
(versus 8\% at $\alpha=1$). Aggregate production (2.30~Mt/yr), wet-biomass
use (6.91--6.96~Mt/yr), and system surplus (708.2, 701.8, and
689.8~M\$/yr for $\alpha=1,3,10$) change by at most 2.6\%, and 94--95\% of the
counties that host facilities are identical across the three cases (58/61 at
$\alpha=3$, 59/63 at $\alpha=10$): the optimal spatial pattern is
feedstock-driven and robust to demand geography. The biochar inherent values
do shift (mean \$324, \$304, and \$476/t for $\alpha=1,3,10$), because they
are nodal prices: at $\alpha=10$ the concentrated demand makes the \$500/t
M-segment bid bind at the demand counties and the inherent value propagates to
the transport-radius limit (\$500/t minus up to \$32/t of delivery cost),
exactly as the market-clearing interpretation predicts.

\subsection{Solver settings}

Gurobi 13.0.2 (academic license), 8 threads. Step 1 MIP: 1\% gap, 1{,}800~s,
MIPFocus=1, NoRelHeurTime=60. Policy sweeps: 0.5\% gap target, 600~s
warm-started from the previous layout (Paradigm~A credit prices
$p_c=25/50$ re-solved with 1{,}800~s limits, improving the incumbents);
sweep points that hit the time limit
are feasible incumbents (reported gaps 0.5--4.7\%, Table~S3). All duals read
from fixed-layout LPs
solved with \emph{dual simplex} (exact vertex shadow prices). The cap's
vertex dual is reported with Gurobi's sign convention for $\le$ constraints in
maximization (the negative of the marginal value): $-\$374$/t at the first
binding cap (near-term), inside the \$351--391/t finite-difference
band of the first binding intervals; the endogenous carbon price is reported as both. The GHG LP
(Step 3) pins
demand quantities to the Step-2 equilibrium (file \texttt{dem\_star\_v2.csv}).
Paradigm C1 is solved both with fixed facilities (LP) and with free facilities
(MIP, C1-free), and Paradigm A and C2 use free-facility MIPs throughout. The
mature-market $p_c=\$100$ point of the Paradigm~A sweep was additionally
re-solved for 1{,}800~s at a 0.1\% gap target
(\texttt{policy\_A\_recert\_v2.csv}); its 500$^\circ$C share improved from
1.75 to 2.79~Mt at a 1.9\% gap.

\begin{table}[htbp]
\centering
\small
\caption{\small Optimality-gap report: free-facility MIPs solve at a 0.5\% gap target
(Step~1 at 1\%); points that hit the time limit are feasible incumbents, so
their reported surplus is a lower bound on the true optimum. v2b runs record
termination status only ($^{\dagger}$), without gap percentages.}
\label{tab:s3}
\begin{tabular}{lcc}
\toprule
experiment (time limit) & near-term & mature-market medium\\
\midrule
Step 1 MIP (1{,}800~s, 1\% target) & OPTIMAL & OPTIMAL\\
Paradigm A sweep (600~s; $p_c=25/50$ at 1{,}800~s) & 0.95--1.9\% & 0.5--3.2\%\\
C1-free (600~s) & 1.1--1.2\% & OPTIMAL\\
C2 sweep (600~s; extended points at 1{,}800~s) & 1.6--4.7\% & OPTIMAL--2.5\%\\
\midrule
\multicolumn{3}{l}{\emph{Quality-differentiated model (v2b)}}\\
Step 1 MIP (1{,}800~s, 1\% target) & TIME\_LIMIT$^{\dagger}$ & OPTIMAL\\
Paradigm A sweep (600~s) & TIME\_LIMIT$^{\dagger}$ & OPTIMAL--TIME\_LIMIT\\
C2 sweep (600~s) & TIME\_LIMIT$^{\dagger}$ & OPTIMAL--TIME\_LIMIT\\
\bottomrule
\end{tabular}

\vspace{2pt}
{\footnotesize $^{\dagger}$Termination status only; gap percentages are not
recorded for the v2b runs.}
\end{table}

\subsection{Eligibility sensitivity: H/C uncertainty and the gate}

The VM0044 gate is applied to the literature-consistent H/C values of
Table~S2. A $\pm$0.1 perturbation around those central values changes
\emph{no} eligibility outcome: every 300$^\circ$C value remains above the gate
($1.00$--$1.40$ minus $0.1$ is still $\ge0.90>0.7$), and every 500$^\circ$C
value remains below it ($0.40$--$0.52$ plus $0.1$ is still
$\le0.62<0.7$). The first flip of a 500$^\circ$C char therefore requires a
$+0.25$ stress, well beyond the published 500$^\circ$C spread of
0.40--0.55, which pushes the agricultural-residue chars (corn stover
$0.49\rightarrow0.74$; soybean/wheat/oats/barley/switchgrass
$\ge0.72$) over the gate while forestry and miscanthus chars
($0.40$--$0.45\rightarrow0.65$--$0.70$) remain eligible. Table~S4 re-solves
Paradigm C2 under that stress. At $p_c=\$50$ nothing changes in either
scenario. At $p_c=\$200$ the canonical full conversion is curtailed: near-term
production falls back to 2.30~Mt of which only 0.31~Mt is 500$^\circ$C
(credits 0.55~Mt, surplus \$716~M versus \$1{,}141~M canonically), and
mature-market production falls back to 2.30~Mt of which 1.59~Mt is
500$^\circ$C (surplus \$1{,}018~M versus \$1{,}501~M). Two conclusions follow.
First, the 300$^\circ$C exclusion is over-determined: the nearest measured
value is 0.9, two tenths above the gate, so no plausible H/C measurement error
re-qualifies torrefaction-like 300$^\circ$C chars. Second, the C2
technology-conversion result is robust within the published 500$^\circ$C H/C
band but degrades if agricultural-residue 500$^\circ$C chars were, contrary to
every measurement we found, above 0.7, which is why VM0044 and IBI treat the
gate as a per-batch measured property rather than a design assumption.

\begin{table}[htbp]
\centering
\small
\caption{\small Paradigm C2 under the $+0.25$ H/C stress (agricultural-residue
500$^\circ$C chars ineligible): surplus, biochar, and 500$^\circ$C share at
$p_c=\$50$ and \$200/tCO$_2$e, versus the canonical values (Table~S2 H/C).}
\label{tab:s4}
\begin{tabular}{llcccc}
\toprule
 & & surplus (M\$) & BC (Mt) & BC500 (Mt) & CC (Mt)\\
\midrule
\multirow{4}{*}{near-term} & canonical, \$50 & 653.3 & 2.30 & 0.00 & 0.00\\
 & stress, \$50 & 653.4 & 2.30 & 0.00 & 0.00\\
 & canonical, \$200 & 1{,}141.1 & 2.25 & 2.25 & 3.64\\
 & stress, \$200 & 715.9 & 2.30 & 0.31 & 0.55\\
\midrule
\multirow{4}{*}{mature-market} & canonical, \$50 & 684.6 & 2.30 & 0.50 & 0.86\\
 & stress, \$50 & 684.5 & 2.30 & 0.50 & 0.86\\
 & canonical, \$200 & 1{,}501.0 & 4.25 & 4.25 & 7.03\\
 & stress, \$200 & 1{,}018.1 & 2.30 & 1.59 & 2.76\\
\bottomrule
\end{tabular}
\end{table}

\subsection{Reproducibility}

One-at-a-time and joint demand sensitivities of the baseline surplus are
reported in the main text (Fig.~\ref{fig:sensitivity}): demand-price shocks
dominate ($\pm$39\% of surplus for $\pm$20\% of price), farmgate and transport
shocks are minor, and a $-40\%$ feedstock-availability shock (a proxy for
field-collection losses) costs 29\% of near-term surplus but only 4\% in the
mature-market scenario. Figures~\ref{fig:s1}--\ref{fig:s3} give the
per-paradigm detail behind the consolidated policy figure: the
baseline-and-credit sweep with credits and net GHG, the tiered-tax and VM0044
technology response, and the cap-and-trade marginal abatement cost curve with
surplus versus the emission cap.

\begin{figure}[htbp]
\centering
\includegraphics[width=0.66\textwidth]{figures/sfig1_policyA.png}
\caption{\small Paradigm A (baseline-and-credit) sweep, near-term: surplus,
production by demand segment, and credits/net GHG as functions of the credit
price. The \$250/t L segment activates at \$25/tCO$_2$e and essentially all
near-term feedstock is mobilized at \$200/tCO$_2$e, tripling surplus.}
\label{fig:s1}
\end{figure}

\begin{figure}[htbp]
\centering
\includegraphics[width=0.72\textwidth]{figures/sfig2_policyC.png}
\caption{\small Paradigm C (tiered tax + VM0044 credit), near-term: technology
choice (I) and surplus (II). The tiered tax leaves the dispatch unchanged, but
the H/C$\le$0.7 gate converts the system to 500$^\circ$C pyrolysis as credit
prices rise.}
\label{fig:s2}
\end{figure}

\begin{figure}[htbp]
\centering
\includegraphics[width=0.78\textwidth]{figures/sfig3_policyB.png}
\caption{\small Paradigm B (cap-and-trade), near-term: (I) marginal abatement
cost curve with the RGGI, EU ETS, and voluntary-market benchmark bands;
(II) system surplus versus the gross-emission cap. The endogenous carbon price
starts at $\approx$\$351/tCO$_2$e and rises to $\approx$\$599/t, above every
observed allowance and voluntary-credit price.}
\label{fig:s3}
\end{figure}

Data generation (\texttt{generate\_data\_v2.py}, \texttt{generate\_data\_v2b.py}),
model (\texttt{01--07\_*.jl}), and figure (\texttt{fig*.py},
\texttt{visualize\_v2.py}) scripts, together with all input data, are available
at \url{https://github.com/YANGHanyu3/Biochar-Supply-Chain}. All optimization models are
implemented in Julia 1.12 with JuMP and solved with Gurobi 13.0.2 (academic
license) on an 8-thread machine; solve times are 1{,}800~s for the Step-1 MIP,
600--1{,}800~s per policy sweep point, and seconds for the fixed-layout LPs.
Figure scripts use matplotlib with the Calibri family at 400 DPI; county
boundaries and the Great Lakes layer are drawn from U.S. Census Bureau and
Natural Earth shapefiles.

\subsection{Proof of Proposition 1}

For the fixed-layout LP, the cap constraint $E^{+}(x)\le\mathrm{CAP}$ is
appended to \eqref{eq:base} with the layout variables $z$ fixed. Strong
duality (Assumption~1) and the differentiability of the LP value function
give, by the standard parametric-LP sensitivity theorem,
$\mu_{\mathrm{CAP}}=\partial\Pi^{*}/\partial\,\mathrm{CAP}$, where
$\mu_{\mathrm{CAP}}$ is the optimal dual variable of the cap constraint.
Gurobi reports the dual of a $\le$ constraint in a maximization with the
opposite sign convention (it reports the reduction in objective per unit of
constraint tightening), i.e., the reported value is
$-\partial\Pi^{*}/\partial\,\mathrm{CAP}$. The endogenous carbon price
(marginal abatement cost) is therefore the negative of the reported dual, and
the finite-difference estimates
$\Delta\Pi^{*}/\Delta\,\mathrm{CAP}$ in Table~S3 reproduce it to numerical
tolerance ($-374$ vs. $\$351$--$391$~\$/t over the first binding intervals,
near-term).

\subsection{Proof of Proposition 2 (two-regime crossover)}

Fix feedstock $f$ and a biochar market price $P$, and consider the two
fixed-yield routes $r\in\{3,5\}$ with dry-biomass yields $\gamma^{r}_f$,
credit rates $\rho^{r}_f$ (tCC/t dry), and chain cost $\omega$ independent of
the route (farmgate is per dry tonne; the route-specific OPEX is already
absorbed in $\omega$; results are unchanged if route-specific costs
$\omega_3\ne\omega_5$ are used, with $\omega(1/\gamma^5-1/\gamma^3)$ replaced by
$\omega_5/\gamma^5-\omega_3/\gamma^3$).

\emph{Demand-bound.} When feedstock is abundant, the active constraint is the
demand bid: the model serves a fixed biochar quantity $Q$, and the supply path
with the higher surplus per tonne of biochar served wins. Per tonne served, the
route-$r$ chain consumes $1/\gamma^{r}_f$ dry tonnes, paying
$\omega/\gamma^{r}_f$ and earning credit revenue
$(\rho^{r}_f/\gamma^{r}_f)\,p_c$, so
\begin{equation}
\Pi^{r}/Q = P - \omega/\gamma^{r}_f + \rho^{r}_f\,p_c/\gamma^{r}_f .
\end{equation}
Subtracting the $r=3$ expression gives
$\Delta^{\mathrm{dem}}_f(p_c)$ as in the statement, which is affine in $p_c$
with slope $\rho^{5}_f/\gamma^{5}_f-\rho^{3}_f/\gamma^{3}_f>0$ (the 500$^\circ$C
credit rate per served tonne exceeds the 300$^\circ$C one; under the VM0044
gate $\rho^3_f=0$ this is automatic). It crosses zero at
$\bar{p}^{\mathrm{dem}}_f$, which is independent of $P$---$P$ cancels in the
difference---and positive whenever the slope is positive. The crossing is
unique because the difference is affine.

\emph{Supply-bound.} When feedstock is scarce, the active constraint is the
feedstock availability: a fixed dry mass $M$ is processed, and the route with
the higher surplus per dry tonne wins. Per dry tonne, route $r$ yields
$\gamma^{r}_f$ t biochar at price $P$ and earns credit revenue
$\rho^{r}_f\,p_c$, so
\begin{equation}
\Pi^{r}/M = \gamma^{r}_f P + \rho^{r}_f p_c - \omega .
\end{equation}
Subtracting gives
$\Delta^{\mathrm{sup}}_f(p_c)=(\gamma^{3}_f-\gamma^{5}_f)P-(\rho^{5}_f-\rho^{3}_f)p_c$,
affine in $p_c$ with \emph{negative} slope, crossing zero at
$\bar{p}^{\mathrm{sup}}_f$, strictly increasing in $P$ (the yield gap
$\gamma^3_f-\gamma^5_f>0$ multiplies $P$), and unique. Note the sign
asymmetry: in the demand-bound regime the credit slope is positive (500$^\circ$C
wins at high $p_c$ because its lower yield per dry tonne concentrates credit
revenue per served tonne), while in the supply-bound regime the credit slope
is negative (the yield loss is a direct revenue loss at price $P$, and credits
must compensate it per dry tonne).

\emph{Numerical values.} Using the pooled feedstock averages of the
LCFS-style baseline-and-credit rates adopted in Paradigm~A
($\bar{\gamma}^{3}=0.352$, $\bar{\gamma}^{5}=0.278$ t/t dry;
$\bar{\rho}^{3}=1.771$, $\bar{\rho}^{5}=1.849$ tCC/t dry, i.e.\ a
credit-rate premium of $\bar{\rho}^{5}-\bar{\rho}^{3}\approx 0.077$ tCC/t
dry; combined chain cost $\omega\approx\$130$/t dry) at the bulk-segment
price $P=\$250$/t:
\begin{equation}
\bar{p}^{\mathrm{dem}}
=\frac{130\,(1/0.278-1/0.352)}{1.849/0.278-1.771/0.352}
\approx\$61/\mathrm{t},\qquad
\bar{p}^{\mathrm{sup}}
=\frac{(0.352-0.278)\,250}{1.849-1.771}
\approx\$239/\mathrm{t}.
\end{equation}
Both values are confirmed by the paradigm-A sweep: the mature-market
(demand-bound) case switches to 500$^\circ$C at $p_c\approx\$61$/t and the
near-term (supply-bound) case stays on 300$^\circ$C through $p_c=\$200$/t with
crossover indicated near $p_c\approx\$239$/t (Fig.~\ref{fig:creditbasis},
which reports the same closed-form crossings with the feedstock averages).
$\square$

\subsection{Proof of Proposition 3 (tax invariance in the demand-constrained regime)}

Write the fixed-layout LP in the compact form
\begin{equation}
\Pi^{*}=\max_{x\in X}\,c^{\top}x,\qquad
X=\{x: Ax\le b,\ x\ge 0\},
\end{equation}
where the demand stack rows are $d\mapsto \mathrm{dem}_d\le \bar{d}_d$ with
objective coefficient $\lambda^{\mathrm{dem}}_d$ in $c$, and the mass-balance
rows carry duals $\pi$. Suppose $x^{*}$ is an optimal basic solution with the
property that the active demand rows are exactly those with
$\pi^{\mathrm{BC}}_{n}<\lambda^{\mathrm{dem}}_{d}$ (strict slack at the margin)
in the terminology of the statement---the demand-constrained regime. Consider
the perturbed problem
\begin{equation}
\Pi^{*}_{\tau}=\max_{x\in X}\,c^{\top}x-\tau\,e^{\top}x,
\quad e\ge 0 \text{ (gross-emission rates per activity unit)}.
\end{equation}
The perturbed objective is $(c-\tau e)^{\top}x$. We show
$\arg\max$ is unchanged. Write $x^{*}=\sum_{\text{basic } j} w_j x_j$; since
$x^{*}$ is an LP vertex with basis $B$, $x_B^{*}=B^{-1}b$ is determined by the
active constraints alone, so any perturbation of the objective---which does not
change the constraint matrix $A$---cannot change the vertex $x_B^{*}$; it can
only change \emph{whether} that vertex is optimal. Optimality persists iff the
reduced costs of the nonbasic activities $j$ remain nonpositive for the
current basis: $c^{\mathrm{red}}_j-\tau e^{\mathrm{red}}_j\le 0$ for all
nonbasic $j$. For the binding-demand rows, the reduced cost of the
corresponding supply path $q$ entering the basis is exactly
$c^{\mathrm{red}}_q=\lambda^{\mathrm{dem}}_{d(q)}-\pi^{\mathrm{BC}}_{n(q)}>0$
(the price gap at the margin); the path is already at its capacity bound and
entering more of it is blocked. All other nonbasic activities have
$c^{\mathrm{red}}_j\le 0$ (they are unprofitable at the current duals), so
they remain unattractive as long as
$\tau\,e^{\mathrm{red}}_j$ does not flip the sign of a positive reduced cost.
Because the only positive reduced costs are the marginal supply paths with
$\lambda^{\mathrm{dem}}-\pi^{\mathrm{BC}}>0$, the critical tax rate is
\begin{equation}
\tau_{\max}=\min_{q\ \text{marginal}}\frac{\lambda^{\mathrm{dem}}_{d(q)}-\pi^{\mathrm{BC}}_{n(q)}}{e^{+}_q}
\end{equation}
as stated; for $\tau<\tau_{\max}$ the optimal basis, and hence the material
flows, technology mix, and cleared quantities, are identical to $x^{*}$, while
the objective decreases by $\tau\,E^{+}(x^{*})$---the tax is a pure
surplus-to-government transfer. Numerically, the marginal price gap is
$\lambda^{\mathrm{M}}-\max_n\pi^{\mathrm{BC}}_n =
500-344.85\approx\$155$/t biochar and the gross-emission intensity of the
marginal supply path is $\approx0.44$~t CO$_2$e/t biochar
($E^{+}=1{,}002.8$~kt over 2.30~Mt served; T1 $31$--$35$~kg/t wet,
T2 $\approx121$~kg/t dry, plus transport), giving
$\tau_{\max}\approx\$350$/t CO$_2$e, above the top tested rate of
\$200/t, consistent with the null dispatch response of C1 and
C1-free (Table~\ref{tab:c1free}). The tiered slope does not matter: a
convex piecewise-linear tax with positive slopes is, on the active paths,
locally linear, so the argument applies tier by tier.
$\square$

\subsection{Nomenclature}

\begin{table}[htbp]
\centering
\small
\caption{Nomenclature: sets, parameters, and variables of the base network
model \eqref{eq:base}. Units: \$ = 2022 USD, t = metric tonne.}
\label{tab:s5}
\begin{tabular}{lll}
\toprule
Type & Notation & Description\\
\midrule
\multirow{9}{*}{Sets} & $\mathcal{N}$ & County nodes ($|\mathcal{N}|=72$) \\
 & $\mathcal{F}$ & Feedstocks ($|\mathcal{F}|=13$: 7 agricultural, 6 forestry) \\
 & $\mathcal{T}$ & Technologies (T1 dehydration; T2 slow pyrolysis at 300/500$^\circ$C; 39 total) \\
 & $\mathcal{P}$ & Products (wet biomass, dry biomass, biochar, carbon credits) \\
 & $\mathcal{S}$, $\mathcal{D}$ & Supply and demand bid rows (farmgate offers, demand segments) \\
 & $\mathcal{A}$ & Transport arcs (radius-limited) \\
 & $\mathcal{K}$ & Facility scale tiers ($k\in\{1,2,3\}$) \\
 & $\mathrm{ref}(t)$ & Reference (input) product of technology $t$ \\
 & $\mathcal{S}_n, \mathcal{D}_n$ & Supply/demand rows located at node $n$ \\
\midrule
\multirow{8}{*}{Parameters} & $\lambda^{\mathrm{dem}}_{d}, \lambda^{\mathrm{sup}}_{s}$ & Demand bid (\$/t) and farmgate offer (\$/t) \\
 & $\omega_t$ & Technology OPEX (\$/t of reference product) \\
 & $\alpha^{\mathrm{var}}_{p}, \alpha^{\mathrm{fix}}_{p}$ & Variable (\$/t-km) and fixed (\$/t) transport cost of product $p$ \\
 & $\kappa_{t,k}$ & Annualized CAPEX of technology $t$ at scale $k$ (\$) \\
 & $\bar{q}_{t,k}$ & Capacity of technology $t$ at scale $k$ (t/yr) \\
 & $\gamma_{t,p}$ & Fixed-yield transformation coefficient of product $p$ by technology $t$ \\
 & $\bar{d}_d, \bar{s}_s$ & Demand and supply capacities (t/yr) \\
 & $\mathrm{AF}$ & Annualization factor ($0.1175$; 10\%, 20 yr) \\
\midrule
\multirow{7}{*}{Variables} & $z_{ntk}\in\{0,\dots,5\}$ & Number of facilities of technology $t$, scale $k$, at node $n$ \\
 & $x_{n,p,t}$ & Net production of product $p$ by technology $t$ at node $n$ (t/yr) \\
 & $f_{ijp}$ & Transport flow of product $p$ from $i$ to $j$ (t/yr) \\
 & $\mathrm{dem}_{d}, \mathrm{sup}_{s}$ & Cleared demand and supply (t/yr) \\
 & $d_{np}, s_{np}$ & Aggregated demand/supply at node $n$ for product $p$ \\
 & $p_{np}$ & Aggregated production at node $n$ for product $p$ \\
 & $\pi_{np}$ & Dual of the mass balance = inherent value of product $p$ at node $n$ (\$/t) \\
\bottomrule
\end{tabular}
\end{table}
